\section{Consent Form for the Human Rating Study}
\label{app:consent}

\noindent\textbf{Study title:} Quality of machine-generated classical Arabic poetry.\\
\textbf{Researcher:} Muhammad Shahbaz, \pending{department and university}; contact \pending{email}.\\
\textbf{Supervisor:} \pending{name}.

\smallskip
\noindent\textbf{What the study is about.} We are studying computer programs that write classical Arabic poems. You
will read about 120 short poems. Some were written by people and some by computer programs; you will not be told which.

\smallskip
\noindent\textbf{What you will do.} For each poem you will give four ratings from 1 (very poor) to 5 (excellent):
overall quality, meaning, fluency, and how the meter sounds when read aloud. This takes about two hours and can be done
in several sittings.

\smallskip
\noindent\textbf{Voluntary participation.} Taking part is voluntary. You may stop at any time without giving a reason
and without any consequence. If you stop, you may ask us to delete the ratings you have given.

\smallskip
\noindent\textbf{Anonymity and data.} Your ratings will be stored under a random code. Your name and contact details
will be kept separately and will not be linked to your ratings in any file we share or publish. Only averages over
raters will be published. The data will be kept for \pending{period} and then deleted.

\smallskip
\noindent\textbf{Risks and benefits.} There are no known risks beyond the time spent. Some poems may contain classical
themes such as war or loss. There is no direct benefit to you; the results will help to build better tools for Arabic
literature and teaching.

\smallskip
\noindent\textbf{Questions.} You can contact the researcher or \pending{the ethics committee} at any time.

\smallskip
\noindent\textbf{Consent.} I have read this information, I have had the chance to ask questions, and I agree to take
part.

\smallskip
\noindent Rater code: \rule{2.5cm}{0.4pt} \quad Signature: \rule{3cm}{0.4pt} \quad Date: \rule{2cm}{0.4pt}

\section{Reproducibility Checklist}
\label{app:repro}
Following Pineau et al.~\cite{pineau2021reproducibility}, we list what is needed to repeat every number in this paper.
\begin{itemize}
\item \textbf{Code:} the baseline notebook \texttt{colab/reproduce\_shaer\_a100.ipynb} (E0--E3) and the research notebook
\texttt{prosody-verifier/prosody\_main.ipynb} (M0--M6), both runnable on Google Colab; the verifier in
\texttt{prosody-verifier/prosody/}; the analysis in \texttt{prosody-verifier/paper/scripts/analysis.py}.
\item \textbf{Data:} the Shaer splits and released poems from the authors' Hugging Face repositories; the pronunciation
lexicon built from the training split and shipped with the code.
\item \textbf{Versions:} every verifier record carries the hashes of the verifier and the lexicon; the notebooks check
their own version against the code.
\item \textbf{Randomness:} generation seed 42; bootstrap seed 7; M4 development subset seed 2026.
\item \textbf{Pre-registration:} the expectations of M3--M6 were committed to the repository before the corresponding
runs, in \texttt{prosody-verifier/validation/}.
\item \textbf{Tests:} the repository's test suite covers the rule tables (textbook equalities), the reading rules, the
rhyme rules, the search, the data builders and the report generators.
\end{itemize}
